\documentclass{article}
\usepackage{graphicx}
\usepackage{parskip}
\usepackage{fullpage}
\usepackage{amsmath}
\usepackage{amsfonts}
\usepackage{makecell}
\usepackage{subfig}
\usepackage{multirow}
\usepackage{caption}
\usepackage{hyperref}
\hypersetup{colorlinks=true, linkcolor=blue, citecolor=blue, urlcolor=cyan,
    pdftitle={Event Contract Mispricing via Options-Implied Probabilities},
    pdfauthor={Tejas Appana}}
\urlstyle{same}
\graphicspath{{figures/}}

\title{Event Contract Mispricing via Options-Implied Probabilities}
\author{Tejas Appana\\ \small Stevens Institute of Technology}
\date{\today}

\begin{document}
\maketitle

\begin{abstract}
\noindent Kalshi, a CFTC-regulated exchange, lists binary contracts on the year-end level of the S\&P 500 in \$200 buckets, and the SPX options market prices the same event far more liquidly. This paper asks whether the two markets price that event consistently, or whether Kalshi prices can be exploited using the information in option prices. Using daily SPXW option chains and Kalshi order books from 2022 to 2024, I extract a risk-neutral density each day with the Breeden--Litzenberger method on a SABR volatility smile and price every Kalshi bucket from it. Three results follow. First, Kalshi stays inside the no-arbitrage band implied by option replication on 97.5\% of bucket-days, so no riskless cross-market arbitrage survives transaction costs. Second, deviations are corrected by Kalshi, not by the options market: Kalshi closes about 4\% of a gap per day, a half-life of roughly one week. Third, the deviations are exploitable by a trader who acts at the daily close: trading Kalshi against the options-implied density earns a pooled excess-return Sharpe ratio of 1.44 (95\% confidence interval 0.36 to 2.49) after fees and the cost of posted collateral, and fills move in the strategy's favor by 0.8 cents per contract net of fees within five days. The edge is short-lived: with a one-day execution delay the Sharpe ratio falls to 0.32 and is no longer significant.

\medskip
\noindent\textbf{Keywords:} prediction markets, event contracts, risk-neutral density, Breeden--Litzenberger, market efficiency, Kalshi\\
\textbf{JEL:} G13, G14
\end{abstract}


\section{Introduction}\label{sec:intro}
Kalshi is a federally regulated exchange on which traders buy and sell ``event contracts'': binary claims that pay one dollar if an event occurs and nothing otherwise, so that a contract's price reads directly as the market's probability of the event \cite{Wolfers2004}. Among its listings are contracts on where the S\&P 500 will close on the last trading day of the year, divided into \$200-wide buckets. The same event is priced, far more liquidly, by the SPX options market: a European option expiring at year-end is a claim on the same closing level, and the full set of option prices across strikes implies a risk-neutral probability distribution for it \cite{Breeden1978}.

Two markets pricing one event invite a simple question, which this paper tests: are they efficient with respect to each other, or does Kalshi misprice the event in a way that the information in option prices can exploit? The options market is treated as the reference. It is orders of magnitude deeper, and the evidence below shows that it does not move toward Kalshi when the two disagree, while Kalshi moves toward it.

The question is answered in three steps, each testing a weaker form of efficiency than the last.
\begin{enumerate}
  \item \textbf{Are the two prices consistent?} Each Kalshi bucket can be replicated, without any model, from option spreads. The replication's bid and ask define a band outside which a riskless cross-market arbitrage would exist (Section~\ref{sec:band}).
  \item \textbf{When they disagree, which market corrects?} An error-correction regression shows which price moves to close a gap, and how quickly (Section~\ref{sec:leadlag}).
  \item \textbf{Is the disagreement exploitable?} A strategy trades Kalshi against the options-implied density, net of fees, the cost of posted collateral, and realistic marking, and is evaluated with statistics that account for the small sample (Section~\ref{sec:strategy}).
\end{enumerate}

The answer is that the two markets are consistent to within transaction costs, but not efficient in the stronger sense: Kalshi lags the options market, and the lag can be traded profitably by a trader who acts at the close. The profit decays within about a day, so it depends on execution speed, and a one-day delay removes most of it.

\subsection{Related Work}
Recovering the risk-neutral distribution of an underlying from the curvature of option prices in strike is due to Breeden and Litzenberger \cite{Breeden1978}; nonparametric and smoothed implementations followed \cite{AitSahalia1995, Melick1997}, and fitting a smooth function to implied volatilities rather than to prices is the standard way to make the second derivative usable \cite{Dumas1998}. Replicating a payoff that depends only on the terminal price from a strip of options is formalized by Carr and Madan \cite{Carr2001}; the bucket replication in Section~\ref{sec:band} is its simplest case. Prediction-market prices are widely read as probabilities \cite{Wolfers2004}, and earlier work has examined arbitrage opportunities within prediction markets themselves \cite{Kildal2012, Cao2013}. This paper instead uses the listed options market on the same underlying event as the reference against which an event-contract market is priced.

\section{Kalshi Event Contracts}\label{sec:kalshi}
Each year-end S\&P 500 market consists of thirteen contracts on \$200-wide buckets (for example, 5600.00 to 5799.99) and two open-ended tail contracts. Exactly one contract pays \$1 at settlement. Kalshi trades continuously, including weekends and after the equity close, through a central limit order book with market makers supplying much of the liquidity. Over the sample period a market order paid a fee of $\lceil 0.035\cdot C\cdot 100\,p(1-p)\rceil$ cents for $C$ contracts at price $p$, and resting limit orders paid none; Kalshi has since doubled the taker rate and introduced a maker fee.

Contracts are fully collateralized: a long position costs its price, and a short position requires posting one dollar minus its price. Over most of the sample that collateral earned nothing, because Kalshi began paying interest on cash and open positions only on 10 October 2024, at a variable 4.05\% annual rate. Tying up collateral therefore has a real cost, which the performance figures in Section~\ref{sec:strategy} charge explicitly.

\section{Data}\label{sec:data}
Two datasets are used: end-of-day SPXW option chains from Refinitiv, obtained through Hanlon Labs at Stevens Institute of Technology, and daily Kalshi candles (closing bid, ask and last-trade price for every contract) from Kalshi's API. The option data contains 344{,}720 quotes from 01/03/2022 to 11/19/2024, with bid and ask prices, vendor implied volatilities, expected dividends, spot and forward levels, and discount factors. The Kalshi data covers the thirteen bucket contracts of the 2022, 2023 and 2024 year-end markets; the two tail contracts are not in the dataset. The Kalshi 2022 market opened on 7 July 2022, and the option feed ends on 19 November 2024, so the three evaluation windows contain 114, 240 and 224 trading days.

Two calendar rules govern which quotes are used. First, an option is retained only if it expires in the same calendar year as its quote date. The feed rolls to the following year's contract about two weeks before the current one settles, and without this rule the last days of each year would be priced from a contract a year from maturity. Second, a new position is opened only on a date on which both markets have a genuine quote. Kalshi trades every day and SPXW options do not; on a day the options market is closed, a gap between a frozen model and a moving Kalshi price no longer indicates that Kalshi is mispriced. Existing positions are still marked and settled on every calendar day.

\subsection{Option Chain Cleaning}
The mid-price is the average of bid and ask. Quotes with only one side are dropped (1{,}892 quotes, almost all ask-only), since they have no mid. A quote whose mid violates the European no-arbitrage bounds --- below intrinsic value on the dividend-adjusted spot, or above the spot for a call and the discounted strike for a put --- is a stale deep in-the-money quote and is dropped (3{,}035 quotes). Quotes more than two standard deviations from that quote date's mean moneyness are excluded, where moneyness is $K/F$ for calls and $2 - K/F$ for puts; the mean and standard deviation are computed within each date, so that no later data affects an earlier one. No bid--ask spread filter is applied.

Implied volatilities are converted from percent to decimal, time to maturity from days to years, and the continuous rate is $r = -\ln(\mathrm{DF})/T$. The spot is adjusted for discrete dividends as $S = S_0 - q\cdot\mathrm{DF}$, where $q$ is the expected dollar dividend before expiry, and the forward is $F = S\,e^{rT}$. Implied volatilities missing from the feed are recovered by inverting Black--Scholes by bisection \cite{Mayhew1995} (4{,}216 quotes). After cleaning, 314{,}626 quotes remain. The effect of each cleaning rule on the results is reported in Appendix~\ref{app:robust}.

\subsection{Kalshi Book Validity}
A Kalshi book is not treated as a price when a side is missing; when it is empty, a bid of at most 2 cents against an ask of at least 98 cents, which arises for buckets the index has left behind; when its spread is at least half the probability range; or on a day when two or more of the mutually exclusive buckets ask 90 cents or more, which cannot all be genuine offers and marks a feed outage. The same rule decides whether a book can generate a trade, mark a position, or enter the market probability. The evidence behind each threshold is in Appendix~\ref{app:book}.

\section{Options-Implied Bucket Probabilities}\label{sec:density}

\subsection{Volatility Smile and Pricing Check}
For each quote date a single smile is fitted to implied volatility as a function of strike using the \emph{out-of-the-money} quotes: puts below the forward, calls above it. By put--call parity the resulting $\sigma(K)$ applies to either option type, and each strike is taken from its more liquid side. In-the-money vendor volatilities are unreliable in this feed --- at some strikes the quoted call volatility differs by as much as nine volatility points from the value implied by the same quote's own mid-price --- which is why they are excluded from the fit. The functional form is SABR \cite{Hagan2002}, using the lognormal expansion with $\beta = 0.5$ and calibrating $\alpha$, $\rho$ and $\nu$ daily. For a single expiry $\beta$ and $\rho$ both act on the slope of the smile and are nearly unidentified jointly, so fixing $\beta$ is standard; results at $\beta = 0$ and $\beta = 1$ are in Appendix~\ref{app:robust}. Beyond the last quoted strike the fitted volatility is held flat. That one fitted smile is used everywhere downstream (Figure~\ref{fig:smoothing}).

Option prices from the smile use Black--Scholes--Merton \cite{Black1973, Merton1973} on the dividend-adjusted spot:
\[ C = S\,N(d_1) - K e^{-rT} N(d_2), \qquad P = K e^{-rT} N(-d_2) - S\,N(-d_1), \]
\[ d_1 = \frac{\ln(S/K) + (r + \sigma^2/2)T}{\sigma \sqrt{T}}, \qquad d_2 = d_1 - \sigma\sqrt{T}. \]
Against observed mids, model prices have a mean absolute error of \$1.99 and an RMSE of \$3.50 across 314{,}624 quotes, including the in-the-money quotes the smile never saw (Figure~\ref{fig:modelmkt}). The errors are smooth in strike rather than random, the signature of a parametric smile that does not bend to every local quote; a curve that did would produce a much noisier second derivative.

\begin{figure}[ht]\centering
  \subfloat[Fitted smile]{\includegraphics[width=0.48\textwidth]{smoothing_example.png}\label{fig:smoothing}}\hfill
  \subfloat[Model versus market prices]{\includegraphics[width=0.48\textwidth]{model_vs_market.png}\label{fig:modelmkt}}
  \caption{SABR smile fitted to out-of-the-money quotes, and the resulting call prices against market mids (06/23/2023).}
\end{figure}

\subsection{Breeden--Litzenberger Density}
Under no-arbitrage, the risk-neutral density of the terminal index is the discounted second derivative of the call price in strike \cite{Breeden1978}:
\[ f_Q(K) = \frac{1}{\mathrm{DF}}\,\frac{\partial^2 C(K)}{\partial K^2} = e^{rT}\,\frac{\partial^2 C(K)}{\partial K^2}. \]
Call prices are evaluated from the fitted smile on a wide strike grid, extended 2{,}000 points below and 1{,}000 above the quoted strikes so that the density has near-full support. A valid density requires $C(K)$ to be convex: its slope $\partial C/\partial K$ must lie in $[-\mathrm{DF}, 0]$ and must not decrease. Inside the quoted strikes the SABR smile satisfies this on all 705 days of the sample. At the last quoted strike it does not, because holding the volatility flat beyond that point puts a kink in the smile, and the kink produces a negative lobe in the second derivative --- a median of 1.8\% of the probability mass per day. The repair is an isotonic projection: the slope is clipped to $[-\mathrm{DF}, 0]$ and replaced by the closest non-decreasing sequence in least squares, computed by the pool-adjacent-violators algorithm, which merges any run in which the slope falls into its average until none remains. The density is the derivative of that slope and is non-negative by construction. On most days the kink lies beyond every listed bucket and the repair leaves bucket probabilities unchanged; on the 220 days on which the quoted strikes stop short of an outer bucket, it moves that bucket's probability by up to 3.7 cents. For the same reason the choice of flat wings has no effect on any result here (Appendix~\ref{app:robust}); it would matter for the open-ended tail contracts, which lie wholly beyond the quotes.

The density is \emph{not} renormalized to integrate to one, and bucket probabilities are \emph{not} rescaled to sum to one: the residual is the probability that the index settles outside every listed bucket, in one of the tails. A bucket's probability is the trapezoidal integral $P([L,U]) = \int_L^U f_Q(K)\,dK$, and a bucket is priced only if it overlaps the quoted strike range. The undiscounted probability is compared directly with the Kalshi price. Strictly, a claim paid at year-end is worth the discounted probability when the collateral behind it earns nothing; the difference averages 0.15 cents per bucket-day, well inside the fee hurdle, and the cost of collateral is charged separately in the performance figures (Figure~\ref{fig:density}).

\subsection{A Skew-Free Baseline}\label{sec:gbm}
To measure what the density contributes, the same smile is collapsed to a single number, the at-the-money volatility $\sigma$ at the strike nearest the forward, and every bucket is priced from a lognormal:
\[ P([L,U]) = \Phi(d(U)) - \Phi(d(L)),\qquad d(K) = \frac{\ln(K/S) - (r - \tfrac12\sigma^2)T}{\sigma\sqrt{T}}. \]
This geometric Brownian motion (GBM) baseline uses the same data and the same smile as the Breeden--Litzenberger density, minus the skew. It is a counterfactual, not a competing forecast: the difference between the two is the value of pricing the skew, and the baseline's bias has a known direction (Section~\ref{sec:gbm2024}).

\begin{figure}[ht]\centering
  \includegraphics[width=0.95\textwidth]{density_example.png}
  \caption{Breeden--Litzenberger density and the bucket probabilities it implies, beside the GBM baseline and the Kalshi mid (06/23/2023). The gaps between the density and Kalshi are the mispricings the strategy trades.}\label{fig:density}
\end{figure}

\section{Are the Two Markets Consistent?}\label{sec:band}
A bucket paying $1[L \le S_T < U]$ is a long digital at $L$ and a short digital at $U$, and each digital can be approximated by a tight spread across the two listed strikes that bracket the edge, scaled by the inverse of the strike gap \cite{Carr2001}. Each edge is built from the out-of-the-money side --- puts below the forward, calls above --- which has the same payoff by put--call parity but avoids the wide spreads of in-the-money quotes; where the two edges use different option types, parity adds a zero-coupon bond leg. Because each leg carries a weight of one over the strike gap, a one-point error in a leg's quote moves the replicated bucket price by ten cents at a ten-point spacing, so the choice of side matters.

Buying the replication at the option asks and selling it at the bids defines a band. If Kalshi's ask were below the replication's bid, or its bid above the replication's ask, after Kalshi's fee, buying one and selling the other would lock in a profit at settlement. Across 4{,}326 bucket-days on which both markets were live and every option leg was two-sided, the Kalshi mid and the replication mid differ by 2.91 cents on average, but the replication costs a median 26 cents round trip in option spreads, against 2 cents on Kalshi. \textbf{Kalshi stays inside the band on 97.5\% of bucket-days.} It sits outside it on 2.5\%, almost all in 2023 (5.2\% of that year's bucket-days, against at most 0.1\% in 2022 and 2024), and outside it by more than the replication's own expected payoff error of 0.8 cents on 0.9\% (Figure~\ref{fig:hedge}). Option commissions are not included, so even these are an upper bound. The two markets are consistent to within the cost of trading between them: there is no riskless cross-market arbitrage.

\begin{figure}[ht]\centering
  \includegraphics[width=0.6\textwidth]{hedging_edge.png}
  \caption{Kalshi price minus the option replication across all bucket-days: the mid-to-mid gap, and the edge that remains after crossing both markets' spreads and paying Kalshi's fee. The x-axis is clipped to $\pm 15$ cents.}\label{fig:hedge}
\end{figure}

\section{Which Market Corrects?}\label{sec:leadlag}
Consistency within a wide band leaves room for the two prices to disagree by several cents. Let $s_t$ be the Kalshi mid minus the Breeden--Litzenberger probability for a bucket on day $t$, over 2{,}344 bucket-days on which the bucket was active in at least one market. An error-correction regression \cite{Engle1987} of each market's subsequent change on $s_t$ shows which one moves to close the gap. The option-implied probability does not move toward Kalshi at all (coefficient $+0.004$, $t = 0.2$). Kalshi does: regressing its change from $t+1$ to $t+2$ on $s_t$ gives $-0.040$ ($t = -2.7$ with standard errors clustered by date, $-3.2$ clustered by date and bucket; 2{,}188 observations over 415 dates), so Kalshi closes about 4\% of a gap per day. The change is measured from $t+1$ rather than $t$ because the Kalshi mid carries bid--ask bounce, which appears in both $s_t$ and the next-day change and makes reversion look faster than it is; the naive next-day regression reads $-0.182$. Instrumenting the spread the same way, with the previous day's value, the median half-life of a gap is 4.9 trading days (interquartile range 3.7 to 7.9). The options market leads and Kalshi follows, over roughly a week.

\section{Is the Mispricing Exploitable?}\label{sec:strategy}

\subsection{Strategy}
On each day on which both markets are live, the strategy buys a bucket when the Breeden--Litzenberger probability exceeds the Kalshi ask by more than the round-trip fee per contract --- two taker fees, one to open and one to close --- and sells it when the probability is below the Kalshi bid by the same margin. A lot is 8 contracts and the account starts with \$200. Exposure is capped at one lot per bucket: a repeated signal on a bucket already held in that direction is ignored, and a signal in the opposite direction closes the position. There is no target, stop or holding limit; a position is held until the opposite signal fires or the contract settles at the true year-end close. A variant that also exits as soon as the mispricing closes was tested and changed nothing (Appendix~\ref{app:robust}). All orders are market orders, since daily candles carry no information about whether a resting order would have filled.

Trades fill at the Kalshi close of the day whose option chain generated the signal. The option quotes are the day's last quotes, probably the 4:15pm close, and the Kalshi daily candle closes at 4:00pm, so this convention can use up to fifteen minutes of information a trader would not yet have had. Section~\ref{sec:timing} bounds its effect. Open positions are marked to the side they would liquidate against --- a long at the bid, a short at the ask --- of the last valid book, and to the model when no valid book has yet been seen.

Returns are measured in excess of the risk-free rate, taken each day from the option chain itself as $r = -\ln(\mathrm{DF})/T$ (3.8\%, 5.5\% and 5.3\% on average over the three windows). The account is treated as two parts: the collateral posted for open positions, which must sit at Kalshi and earns Kalshi's rate (zero before October 2024), and idle cash, which is assumed to earn the risk-free rate elsewhere. The daily excess return is therefore the trading profit plus Kalshi's interest on the collateral minus the risk-free rate on it; carrying the collateral costs between 0.6\% and 1.9\% a year. Returns are sampled on trading days over the window in which both markets were live and annualized with 252 days; alpha and beta against SPX use Newey--West standard errors \cite{NeweyWest1987} with the lag set by the Newey--West rule \cite{NeweyWest1994}. Sharpe intervals are 95\% stationary block bootstraps \cite{Politis1994} with 10{,}000 resamples and a mean block of ten days.

\subsection{Results}
A single year of 114 to 240 trading days is too short for a Sharpe interval to exclude zero even at a ratio near 2, so the headline test pools the three years' daily excess returns. \textbf{The pooled Sharpe ratio is 1.44, with a 95\% interval of 0.36 to 2.49; 99.5\% of bootstrap resamples are positive} (Table~\ref{tab:pooled}). The same test on the GBM baseline gives 0.70, with an interval that includes zero. The result does not depend on the modelling choices: across every variant in Appendix~\ref{app:robust} --- cleaning rules, SABR $\beta$, smile form, lot size, Kalshi's current doubled fee (pooled 1.16), the exit rule and the marking of unquoted buckets --- the pooled interval stays above zero. The only change that removes it is delaying execution (Section~\ref{sec:timing}).

\begin{center}
\begin{tabular}{l|c|c|c}
    & Same close & Model one day old & Model two days old \\
    \hline
    Breeden--Litzenberger & \textbf{1.44} [0.36,\,2.49] & 0.32 [$-$0.70,\,1.34] & 0.03 [$-$0.97,\,1.05] \\
    GBM baseline & 0.70 [$-$0.58,\,1.88] & $-$1.02 [$-$2.09,\,0.11] & $-$1.12 [$-$2.22,\,0.01] \\
\end{tabular}
\captionof{table}{Pooled 2022--2024 excess-return Sharpe ratio with 95\% block-bootstrap interval (575 trading days), by execution timing.}\label{tab:pooled}
\end{center}

The trades themselves carry the same evidence. After each fill, the Kalshi mid moves in the trade's favor by 0.66 cents per contract the next day ($t = 2.3$, clustered by fill date) and 1.21 cents after five days ($t = 3.1$), which is 0.83 cents net of the fill's own fee ($t = 2.1$), and then stops. By year (Table~\ref{tab:perf}), the strategy earns a positive excess return in every year, with Sharpe ratios of 1.98, 1.90 and 0.55; 2022 covers under half a year and deserves the least weight.

\begin{center}
\small
\begin{tabular}{c|c|c|c|c|c|c|c|c}
    Year & Excess return & Ann.\ vol & Sharpe & Sharpe 95\% CI & Max DD & Total$^\dagger$ & Beta & Fills \\
    \hline
    2022$^\ddagger$ & 13.16\% & 6.66\% & 1.98 & [0.20,\,3.90] & $-$1.87\% & 9.82\% & 0.02 & 98 \\
    2023 & 6.54\% & 3.45\% & 1.90 & [$-$0.28,\,3.66] & $-$3.14\% & 9.20\% & 0.03 & 81 \\
    2024$^\S$ & 1.65\% & 3.00\% & 0.55 & [$-$1.05,\,2.04] & $-$1.56\% & 3.47\% & 0.02 & 40 \\
    \hline
    \multicolumn{9}{l}{\footnotesize GBM baseline Sharpe: 2.31, 1.08, $-$1.98}\\
\end{tabular}
\captionof{table}{Breeden--Litzenberger strategy by year (one lot per bucket, both sides, same-close execution). Excess return and volatility are annualized on trading days in excess of the option-implied rate, with posted collateral charged its cost of capital; maximum drawdown is on the full calendar path. $^\dagger$Nominal end-to-end return of the Kalshi account including settlement of held buckets. $^\ddagger$From 7 July 2022 (114 trading days). $^\S$Return path ends 19 November 2024 with the option feed; positions settle at the true year-end close, and SPX finished above every listed bucket.}\label{tab:perf}
\end{center}

The strategy is close to market-neutral: betas are at most 0.03, alpha against SPX has Newey--West $t$-statistics between 0.4 and 1.8, and drawdowns are at most 3.1\%. Both sides contribute. The sell side is positive in every year (Sharpe 0.31, 0.56 and 0.38 run alone), and the buy side is positive in 2022 and 2023 and negative in 2024 (0.27, 0.93 and $-$1.19). Figure~\ref{fig:strat_bench} shows the account against SPX.

\begin{figure}[ht]\centering
  \includegraphics[width=0.98\textwidth]{strategy_vs_benchmark.png}
  \caption{Breeden--Litzenberger strategy (nominal account value) against SPX normalized to the same \$200, by year.}\label{fig:strat_bench}
\end{figure}

\subsection{Execution Timing}\label{sec:timing}
The claim that the mispricing is exploitable rests on acting at the close. Delaying every trade removes any possible lookahead: each day's signal is computed from the previous day's option chain and re-tested against the Kalshi book it would actually fill against. With a one-day delay the pooled Sharpe ratio falls to 0.32 and with two days to 0.03, neither significant (Table~\ref{tab:pooled}). The markouts show why. A delayed fill moves \emph{against} the trade by 0.62 cents the next day ($t = -3.4$): most of a mispricing is corrected within a day, so a trade placed a day late enters after the correction has begun.

The one-day delay is a conservative bound, not the relevant case. The lookahead in the same-close convention is fifteen minutes, not a day, and a trader could compute the density from option quotes shortly before the Kalshi close. Where the truth lies between the two columns of Table~\ref{tab:pooled} can be measured with intraday data; the finding here is that the mispricing is exploitable only by a trader who acts on the day it appears.

\subsection{Risk}
Collateral never exceeds 45\% of the account in any year, so the results use no hidden leverage. Because the buckets are mutually exclusive, a settle-now stress --- supposing on each date that the year settled in the worst bucket for the open book --- is bounded by the one-lot cap: its worst value is $-3.2\%$, $-3.1\%$ and $-5.5\%$ of capital. The exit rule matters. Closing a position when the opposite signal fires, rather than holding everything to settlement, added \$49.32 across the three years on the \$200 base and turned 2022 from $-7.6\%$ into $+9.8\%$; it acts as a model-based stop, though 16 of the 36 positions it closed at a loss would have ended profitable if held.

Two measurement effects qualify the daily Sharpe ratio. Daily returns are autocorrelated, so the $\sqrt{252}$ scaling is only approximate \cite{Lo2002}; Sharpe ratios from non-overlapping 20-day returns have the same sign in every year (3.57, 1.17 and 0.58). And late in each year the buckets the index has left stop being quoted, so a position there is carried at its last valid book: 18\%, 36\% and 66\% of position-days in the three years. A mark that does not move understates volatility. Marking those days to the model instead gives daily Sharpe ratios of 2.15, 0.99 and 0.44, which errs the other way, since the position then switches between a liquidation price and a model probability as the book comes and goes; the 20-day figures barely change (3.77, 1.09 and 0.44). The 2023 daily Sharpe of 1.90 is best read beside its 20-day value of about 1.1.

\subsection{What the Density Adds}\label{sec:gbm2024}
The GBM baseline trades the same smile without its skew and earns a pooled Sharpe ratio of 0.70, against 1.44. The gap is largest in 2024, when the baseline loses money, and the reason is a bias with a known direction. On 11/19/2024 the fitted smile ran from 28.1\% at a strike of 5000 to 14.3\% at 5800, against an at-the-money volatility of 11.8\%. A lognormal that narrow concentrates probability near spot and overprices buckets moderately below it: the 5600--5799.99 bucket, 3.6\% below spot, was worth 20.0 cents under the baseline against 10.5 under the density. The bias is systematic. Across all three years the baseline exceeds the density by 4.7 cents on average for buckets 2 to 5\% below spot and by 4.3 cents for buckets 5 to 10\% below, and understates buckets above spot by about 3 cents. What made 2024 different is where the buckets sat: SPX rallied from 4743 to 5917, so 82\% of 2024 bucket-days were below spot, against 58 to 60\% in the other years. The baseline read its overpricing of those buckets as cheapness; 14 of its 16 purchases were below-spot buckets, and they lost \$4.72 of its \$5.20 net trading loss when SPX settled above every bucket. Pricing the skew is what separates a profitable signal from a losing one.

\section{Conclusion}\label{sec:conclusion}
Kalshi's year-end S\&P 500 contracts and the SPX options market price the same event consistently to within transaction costs: Kalshi stays inside the no-arbitrage band implied by option replication on 97.5\% of bucket-days, and the replication costs an order of magnitude more to trade than the gap it would capture. They are not efficient in the stronger sense. When the two disagree, Kalshi does the correcting, at about 4\% of a gap per day, and a strategy that trades Kalshi against the options-implied Breeden--Litzenberger density earns a pooled excess-return Sharpe ratio of 1.44 after fees and the cost of collateral, with an interval that excludes zero. Kalshi is therefore exploitably mispriced relative to the options market, for a trader who acts at the close.

Three qualifications bound the claim. The edge is short-lived: with a one-day execution delay it is no longer significant, so it belongs to a trader who can act on the day a mispricing appears. The sample holds three years and three settlement outcomes, and 2022 covers under half a year. And the scale is small: one lot per bucket on a \$200 base, in order books thin enough that the relationship between the model and realizable prices is loose.

Three extensions follow directly. Hourly Kalshi candles would allow trading after the option close on the same day, measuring the edge without any lookahead instead of bounding it. The two open-ended tail contracts, which the density already prices, would complete each market and include the contract on which 2024 actually settled. And additional years, or Kalshi's shorter-dated S\&P 500 markets, would multiply the number of settlement outcomes, which is what the statistical claims here need most.

\section*{Acknowledgments}
I thank Hanlon Labs at Stevens Institute of Technology for access to the option data, and Adam, who worked with me on the preliminary hypothesis. All errors are my own. Code and documentation are available at \url{https://github.com/tej100/kalshi-stat-arb}.

\bibliographystyle{plain}
\bibliography{ref}

\newpage
\appendix
\section{Robustness}\label{app:robust}
Each row changes one assumption against the headline configuration and reports the Breeden--Litzenberger strategy's Sharpe ratio by year, the pooled Sharpe ratio with its 95\% block-bootstrap interval, and the option pricing error. None of these variants was used to choose the headline configuration.

\begin{center}
\small
\begin{tabular}{llrrrcr}
& Variant & 2022 & 2023 & 2024 & Pooled [95\% CI] & MAE \\
\hline
Headline & as in the paper & 1.98 & 1.90 & 0.55 & 1.44 [0.36,\,2.49] & \$1.99 \\
\hline
Cleaning & spread filter 30\% of mid & 1.89 & 1.90 & 0.68 & 1.45 [0.38,\,2.49] & \$1.85 \\
 & moneyness cut 3 sigma & 2.00 & 1.68 & 0.61 & 1.39 [0.31,\,2.47] & \$2.08 \\
 & no moneyness cut & 1.91 & 1.66 & 0.62 & 1.36 [0.30,\,2.41] & \$2.09 \\
\hline
Smile & SABR $\beta$ = 0 & 1.97 & 1.99 & 0.50 & 1.46 [0.40,\,2.48] & \$1.93 \\
 & SABR $\beta$ = 1 & 1.57 & 1.68 & 0.55 & 1.23 [0.16,\,2.28] & \$2.22 \\
 & SABR wings (no flat extrapolation) & 1.98 & 1.90 & 0.55 & 1.44 [0.36,\,2.49] & \$1.99 \\
 & SVI instead of SABR & 2.06 & 2.47 & 0.38 & 1.64 [0.53,\,2.73] & \$1.90 \\
\hline
Sizing & lot of 1 contract (capital \$25) & 1.32 & 1.61 & 0.34 & 1.04 [0.03,\,2.05] & \$1.99 \\
 & lot of 4 contracts (capital \$100) & 1.87 & 1.72 & 0.47 & 1.31 [0.23,\,2.40] & \$1.99 \\
 & lot of 16 contracts (capital \$400) & 2.02 & 1.93 & 0.57 & 1.47 [0.40,\,2.51] & \$1.99 \\
\hline
Costs & current Kalshi fee (0.07) & 1.51 & 1.66 & 0.40 & 1.16 [0.10,\,2.21] & \$1.99 \\
\hline
Exit rule & convergence exit & 1.91 & 1.99 & 0.59 & 1.46 [0.43,\,2.49] & \$1.99 \\
\hline
Execution & model 1 day old & 0.67 & $-$0.13 & 0.52 & 0.32 [$-$0.70,\,1.34] & \$1.99 \\
 & model 2 days old & 0.00 & $-$0.22 & 0.37 & 0.03 [$-$0.97,\,1.05] & \$1.99 \\
\hline
Marking & stale marks to model & 2.15 & 0.99 & 0.44 & 1.04 [0.19,\,1.90] & \$1.99 \\
\end{tabular}

\captionof{table}{Sensitivity of the Breeden--Litzenberger strategy to each modelling choice. Lot-size variants scale capital with the lot. The convergence exit also closes a position as soon as the bid (for a long) or ask (for a short) reaches the model probability.}\label{tab:robust}
\end{center}

Seven smile functional forms were compared on the same quotes and the same density construction: least-squares and cubic smoothing splines \cite{Dierckx1995}, a degree-4 polynomial, PCHIP, LOWESS, SVI \cite{Gatheral2014} and SABR. SABR and SVI produced the smoothest densities with the fewest spurious modes (roughness 2.17 and 2.19, against 2.41 to 31.2 for the others) and the smallest share of negative density (0.002); SABR calibrated on every date. The strategy results for SVI, for the average of the SABR and SVI probabilities, and for trading only where both agree lie within 0.61 of SABR in every year, against intervals about 3.6 wide.

\section{Kalshi Book Validity}\label{app:book}
Besides books with a missing bid or ask, three quote artifacts in the Kalshi data would otherwise be read as prices. Each is excluded, and a position whose book is excluded is marked to the last valid book, or to the model if none has yet been seen.
\begin{itemize}
  \item \textbf{Empty books.} Once the index has moved far past a bucket, no participant offers it, and its book reads a zero bid against a 100-cent ask. Treating that ask as a liquidation price marks a short position at a near-total loss that no trade could have realized. Books with a bid of at most 2 cents and an ask of at least 98 cents are excluded.
  \item \textbf{Stray offers.} Isolated offers on buckets the index had long passed quote asks of 85 to 86 cents against zero bids. Every legitimate spread in the sample is at most 44 cents, so a book whose spread is at least 50 cents is excluded; any cutoff from 45 to 84 cents gives identical results. Without the rule a single such offer produces a \$6.80 one-day loss, the whole of the 2024 maximum drawdown.
  \item \textbf{Feed outages.} On 16 November 2024 all thirteen buckets read a zero bid against a 100-cent ask while their last-trade prices stayed between 1 and 13 cents, and the outage decayed through 21 November with asks of 97 to 99 cents. Because exactly one bucket can pay, two or more buckets asking 90 cents or more cannot both be genuine; such days are excluded for every bucket.
\end{itemize}

\end{document}
